\begin{afterword}
\summaryhead

The post tells two short stories about Ougi, a master of rationality, in the style of Zen koans.

In the first, a friend tells a novice that he has fallen into a cult. The novice asks Ougi for reassurance. Ougi answers with a riddle about a hammer, then promises plain arguments if the novice will wear a silly cowboy hat. The novice puts it on. Ougi then explains: the novice came for reassurance, not out of curiosity, and his need for it made him willing to obey; an evil master could have taken his money. The worth of an understanding shows when it is used on real questions. The novice later becomes a master himself and forbids his students to quote him: ``Use the techniques and do not mention them.''

In the second, a novice complains that the dojo's robes and hoods look cultish and irrational. Ougi hands him a wizard's hat, since someone so worried about how clothing interacts with probability theory must need a special hat. That novice later discusses rationality only in a clown suit.

\responsehead

The first koan makes a real point, and makes it by demonstration. A person anxious to be told that their group is not a cult will obey an arbitrary order to get the answer, and that readiness, not the doubt, is what a manipulator uses. The master says as much himself (``If I had been an evil man''). The remedy offered is curiosity and use: test the method on real questions, and stop watching for who is laughing.

The second koan answers a different worry from the one the novice raises. His complaint is that the dojo is ``a little cultish,'' with robes and hoods like the Freemasons. The master treats this as a claim that clothing affects reasoning, and mocks it. But nine days earlier, ``Every Cause Wants To Be A Cult'' traced cultishness to ordinary in-group feeling, not to faulty reasoning. Whether shared dress feeds that feeling is a fair question, whatever dress does to probability theory, and the koan does not take it up.

The second koan's ending can be read either way: the novice in the clown suit may have learned that dress does not matter, or learned it too well. The author states the first koan's lesson plainly nine days later, in ``Cultish Countercultishness,'' which repeats the image of the swordsman and points back to these koans. A reader who wants the argument will find it there.

In short: a first koan with a real point about the need for reassurance, and a second that answers a worry about in-group feeling as if it were a worry about probability theory.
\end{afterword}
